The article predicts at one moment --- case creation --- from the attributes
recorded there, and Section~\ref{sec:lim} states plainly that prefix length,
prefix bucketing and sequence encoding are therefore not axes of
$\mathcal{S}$ and cannot be. That limitation stands. What it leaves open is
the question a reader from predictive process monitoring will ask next:
whether the paper's central claim --- that the moment of prediction is an
\emph{argument} of the estimand rather than a preliminary to it --- is a
property of that literature's own axis as well, or an artefact of predicting
at creation. This section answers that on one log.

\paragraph{The construction, declared before it was run}
The log is \prefixLog, the case study's own, and the register is the affected
configuration item. At prefix $k$ the analyst stands after the case's $k$-th
assignment event. Three things follow from $k$ and are not chosen
separately. The \textbf{population} is the cases still at risk: more than $k$
events, and no change of assignment group in the first $k$, because a case
whose outcome is already determined by $k$ is not a prediction problem. The
\textbf{target} is whether a change of group occurs \emph{after} event $k$,
which is the registered handover rule of Section~\ref{sec:design} restricted
to the remaining suffix, so nothing already observed is predicted. The
\textbf{baseline} is the intake block plus what the prefix has revealed ---
the group the case sits in now and how many distinct groups it has seen ---
which is the prefix-aware baseline that literature builds, and giving the
prefix to the baseline and not to the register is the comparison that can
embarrass the register. The split is the same fixed temporal holdout under
the same declared tie-break, and every interval is the pointwise basic
construction under the block-weighted bootstrap at \nPrefixDraws\ draws.
$k = 0$ is case creation, which is the prediction point every other result in
this paper occupies.

\begin{table}[t]
\centering
\small
\resizebox{\ifdim\width>\linewidth\linewidth\else\width\fi}{!}{%%
\begin{tabular}{llrlrrrlr}
\toprule
prefix & n & prevalence & register levels & base auc & with auc & V & interval & resolved \\
\midrule
0 & 46606 & 0.945 & 3019 & 0.641 & 0.898 & 0.257 & [+0.233, +0.281] & True \\
1 & 46605 & 0.945 & 3019 & 0.733 & 0.900 & 0.166 & [+0.145, +0.203] & True \\
2 & 11080 & 0.722 & 1567 & 0.711 & 0.775 & 0.064 & [+0.047, +0.088] & True \\
3 & 6583 & 0.612 & 1224 & 0.799 & 0.816 & 0.017 & [-0.025, +0.039] & False \\
5 & 3005 & 0.373 & 700 & 0.743 & 0.787 & 0.044 & [+0.021, +0.087] & True \\
8 & 857 & 0.550 & 341 & 0.672 & 0.741 & 0.069 & [-0.001, +0.126] & False \\
\bottomrule
\end{tabular}%
}
\caption{A PREFIX AXIS ON ONE LOG. At prefix $k$ the analyst stands after the case's $k$-th assignment event: the population is the cases still at risk --- more than $k$ events and no group change yet --- the target is whether a change occurs AFTER event $k$, and the baseline is the intake block plus what the prefix has revealed. $k = 0$ is case creation, which is the prediction point every other result in this paper uses, and it is the anchor the rest are read against. Every interval is the pointwise basic construction under the block-weighted bootstrap at \nPrefixDraws\ draws.}
\label{tab:prefix}
\end{table}


\paragraph{What it shows}
\textbf{The register's worth falls sharply away from case creation, and
what it does after that is not measured well enough to describe.} At case
creation it is \VPrefixZero\ [\VPrefixZeroLo, \VPrefixZeroHi] AUC over the
intake block on \nPrefixZero\ cases; two events later it is \VPrefixTwo\
[\VPrefixTwoLo, \VPrefixTwoHi] on \nPrefixTwo. That fall is resolved: the two
intervals do not overlap. \textbf{Beyond it the sequence is not monotone} ---
at the deepest prefix measured the increment is \VPrefixEight\
[\VPrefixEightLo, \VPrefixEightHi], nominally \prefixRetainedPct\ of its
creation-time value but on an interval containing zero, and higher than the
point estimate at three shallower prefixes. Across
the \nPrefixLevels\ prefixes it ranges over \prefixVmin\ to \prefixVmax, and
on \nPrefixUnresolved\ of them the interval does not exclude zero, so no
shape should be read into the tail of this axis --- only the fall at its
head, which is determined, and large.

\textbf{Two of the three mechanisms are visible in the table.} The baseline
rises as the prefix reveals the routing history, from \basePrefixZero\ at
creation; the population shrinks and its prevalence falls as the cases that
were going to be handed over early leave it; and the register's own
cardinality falls with the population. None of this is drift in the register:
it is the same field, on the same estate, worth a different amount because the
question changed underneath it.

\paragraph{What this does and does not license}
It does not make this a predictive-process-monitoring paper, and no claim
about that literature's benchmarks is made here: one log, one register, one
encoding, no prefix bucketing and no sequence encoder is a pilot and is
offered as one. What it does establish is that the prediction point is an
argument of the estimand on that literature's own axis and not only on the
service-desk one of Section~\ref{sec:twodecision}, and that case creation ---
the point this paper occupies throughout --- is the \emph{most favourable}
point for a creation-time register. A study that reported this register's
worth from a prefix-based benchmark and a study that reported it from case
creation would disagree by a factor of \prefixDropFactor\ \emph{in point
estimates}, and neither would be wrong. The qualifier is not decoration: that
factor's denominator is one of the prefixes whose interval contains zero.
